\documentclass[12pt,a4paper]{article}

% ============================================================================
% PACKAGES
% ============================================================================
\usepackage[utf8]{inputenc}
\usepackage[T1]{fontenc}
\usepackage{lmodern}
\usepackage[margin=1in]{geometry}
\usepackage{setspace}
\usepackage{graphicx}
\usepackage{amsmath,amssymb}
\usepackage{booktabs}
\usepackage{hyperref}
\usepackage{cleveref}
\usepackage{natbib}
\usepackage{enumitem}
\usepackage{xcolor}

\hypersetup{
    colorlinks=true,
    linkcolor=blue!70!black,
    citecolor=blue!70!black,
    urlcolor=blue!70!black
}

\onehalfspacing

% ============================================================================
% DOCUMENT
% ============================================================================
\begin{document}

\section*{Introduction}

\subsection{Background and Motivation}

The rapid evolution of Artificial Intelligence has positioned Automatic Speech Recognition (ASR) as a cornerstone of modern human-computer interaction, enabling applications that span virtual assistants, real-time transcription services, automated customer support, and accessibility technologies for persons with disabilities. Over the past decade, ASR systems for high-resource languages such as English, Mandarin, and Spanish have achieved near-human levels of transcription accuracy, driven by the convergence of three factors: the availability of tens of thousands of hours of annotated speech data, the development of increasingly powerful neural architectures, and access to large-scale computational infrastructure \citep{baevski2020wav2vec, radford2023robust}. However, this remarkable progress has been profoundly uneven in its distribution across the world's linguistic landscape. Of the approximately 7,000 languages spoken globally, fewer than 100 possess the annotated speech corpora, standardized orthographies, and computational linguistic resources necessary to train competitive ASR systems \citep{joshi2020state}. This asymmetry constitutes a significant digital divide, one that excludes billions of speakers from the benefits of voice-enabled technology and reinforces existing patterns of linguistic marginalization.

The African continent exemplifies this disparity with particular acuity. Africa is home to over 2,000 languages distributed across four major language families---Niger-Congo, Afro-Asiatic, Nilo-Saharan, and Khoisan---yet the overwhelming majority of these languages remain entirely absent from the global speech technology ecosystem \citep{adelani2022masakhaner}. Even languages with millions of speakers, such as Yoruba, Igbo, Amharic, and the various Nilotic languages of East Africa, are classified as ``low-resource'' or ``under-resourced'' in the computational linguistics literature, lacking the structured datasets, annotated corpora, and established processing pipelines that underpin modern ASR development \citep{magueresse2020low}. This exclusion carries consequences that extend far beyond technological inconvenience. In contexts where oral communication remains the primary mode of information exchange---as is the case across much of rural Africa---the absence of speech technology in indigenous languages limits the deployment of critical services in education, healthcare, agriculture, and governance. Voice-based interfaces could enable illiterate or semi-literate populations to access digital information systems, yet this potential remains unrealized for the vast majority of African language communities.

Kalenjin, a Southern Nilotic language spoken primarily in Kenya's Rift Valley region, serves as a representative and instructive case study for these challenges. With an estimated 5 to 6 million speakers distributed across several closely related sub-varieties---including Nandi, Kipsigis, Tugen, Keiyo, Marakwet, Pokot, and Sabaot---Kalenjin constitutes one of the larger indigenous language groups in East Africa \citep{toweett1979study}. Despite this substantial speaker population, Kalenjin occupies a paradoxical position: it is widely spoken in daily life, used in local media broadcasts, and serves as a language of cultural identity for the Kalenjin community, yet it remains severely under-resourced from a computational perspective. No large-scale annotated speech corpus existed for Kalenjin prior to its inclusion in the Mozilla Common Voice project, and the language lacks the extensive text corpora, pronunciation dictionaries, and morphological analyzers that facilitate ASR development for better-resourced languages.

\subsection{Linguistic Characteristics of Kalenjin}

The linguistic properties of Kalenjin present specific challenges for automatic speech recognition that merit detailed consideration, as they inform the design decisions throughout the preprocessing and modeling pipeline developed in this work.

Kalenjin belongs to the Southern Nilotic branch of the Nilo-Saharan language family, a classification that places it within a typological context distinct from the Bantu languages that have received comparatively greater attention in African language technology research \citep{rottland1982southern}. As an agglutinative language, Kalenjin constructs complex words through the systematic concatenation of morphemes, with a single word potentially encoding information about tense, aspect, mood, person, number, and case through affixation. This morphological richness means that the effective vocabulary of the language is substantially larger than a simple word list would suggest, posing challenges for word-level language modeling and increasing the importance of character-level or subword-level modeling approaches in ASR.

Kalenjin exhibits tonal properties, wherein pitch variations at the syllabic level carry lexical and grammatical significance. Subtle differences in fundamental frequency can distinguish otherwise identical segmental sequences, producing distinct meanings from phonetically similar utterances. This tonal dimension necessitates high-fidelity acoustic modeling capable of capturing suprasegmental features that extend beyond the spectral characteristics typically emphasized in ASR systems designed for non-tonal languages. The interaction between tone and the segmental phonology of Kalenjin creates a complex acoustic space that the model must learn to navigate, a task made more difficult by the limited training data available.

The phonological inventory of Kalenjin includes several sounds that are typologically uncommon in the high-resource languages that dominate ASR pre-training corpora. The velar nasal /\ng/, orthographically represented as \textit{ng'} with a distinctive apostrophe marker, is phonemically contrastive with the sequence /ng/ and appears with high frequency in Kalenjin text. The preservation of this orthographic distinction is critical for accurate transcription and informed the text normalization decisions described in the methodology. Additionally, Kalenjin employs a vowel harmony system based on the Advanced Tongue Root (ATR) feature, where vowels within a word must agree in their ATR specification, a constraint that influences both phonetic realization and morphological processes.

The dialectal variation across Kalenjin sub-varieties introduces further complexity. While the sub-varieties are largely mutually intelligible, they exhibit systematic differences in phonological rules, lexical items, and morphological patterns. A crowd-sourced speech corpus such as Mozilla Common Voice inevitably contains contributions from speakers of multiple sub-varieties, introducing acoustic and linguistic variation that the ASR system must accommodate without explicit dialect labels or dialect-specific modeling.

Finally, the absence of a fully standardized orthography for Kalenjin means that the same word may appear in multiple written forms across different sources. While a broadly accepted orthographic convention exists based on the Latin alphabet, inconsistencies in the representation of certain sounds---particularly the treatment of aspirated consonants, the notation of vowel length, and the use of apostrophes---create challenges for text normalization and vocabulary construction. These orthographic inconsistencies must be resolved systematically to produce a clean, consistent character set suitable for Connectionist Temporal Classification (CTC) decoding.

\subsection{Evolution of ASR Architectures}

The development of automatic speech recognition has undergone several paradigmatic shifts over the past four decades, each expanding the range of languages and acoustic conditions that ASR systems can handle. Understanding this trajectory is essential for contextualizing the architectural choices made in this work.

The dominant paradigm from the 1980s through the early 2010s was the Hidden Markov Model--Gaussian Mixture Model (HMM-GMM) framework, in which speech was modeled as a sequence of discrete states with Gaussian-distributed acoustic observations \citep{rabiner1989tutorial}. While effective for controlled environments and well-resourced languages, HMM-GMM systems required extensive feature engineering, hand-crafted pronunciation dictionaries, and separate acoustic, language, and pronunciation models that were trained independently and combined during decoding. The resource requirements of this pipeline---particularly the need for phonetic transcriptions and expert linguistic knowledge---made it largely inaccessible for under-resourced languages.

The introduction of deep neural networks (DNNs) as replacements for Gaussian mixture models in the acoustic modeling component marked the first major shift, with hybrid DNN-HMM systems achieving substantial accuracy improvements \citep{hinton2012deep}. However, these systems retained the modular pipeline architecture and its associated resource requirements. The subsequent development of end-to-end (E2E) architectures---including Connectionist Temporal Classification \citep{graves2006connectionist}, attention-based encoder-decoder models \citep{chan2016listen}, and the RNN-Transducer \citep{graves2012sequence}---eliminated the need for separate components by directly mapping acoustic input sequences to character or subword output sequences. CTC, in particular, proved well-suited for low-resource scenarios due to its ability to handle variable-length alignment between input and output sequences without requiring frame-level phonetic annotations.

The most transformative development for low-resource ASR has been the emergence of self-supervised learning for speech representations. Wav2Vec2 \citep{baevski2020wav2vec} introduced a framework in which a model learns universal speech representations from large quantities of unlabeled audio through a contrastive learning objective applied to masked latent representations. The model architecture consists of a convolutional neural network (CNN) feature encoder that processes raw audio waveforms into latent representations, followed by a Transformer encoder that captures long-range contextual dependencies through self-attention. By pre-training on thousands of hours of unlabeled speech, the model acquires robust acoustic representations that can be fine-tuned for specific downstream tasks---including ASR---with relatively small amounts of labeled data.

The cross-lingual extension of this approach, XLS-R \citep{babu2022xls}, scaled self-supervised pre-training to 436,000 hours of speech data spanning 128 languages, demonstrating that multilingual pre-training produces representations that transfer effectively across languages, including those not seen during pre-training. This cross-lingual transfer capability is particularly significant for under-resourced languages, as it allows researchers to leverage the acoustic knowledge encoded in multilingual representations to bootstrap ASR systems with minimal target-language data. The XLS-R-300M variant, with 300 million parameters distributed across 7 CNN layers and 24 Transformer layers, provides a powerful foundation model that balances representational capacity with fine-tuning feasibility for resource-constrained settings.

\subsection{The Low-Resource ASR Challenge}

Building ASR systems for under-resourced languages involves a constellation of interrelated challenges that extend beyond the simple scarcity of training data. While data volume is undeniably a primary constraint---high-resource ASR systems typically train on 10,000 to 100,000 hours of transcribed speech, compared to the fewer than 20 hours available for Kalenjin---the problem is fundamentally multidimensional.

\textbf{Data quality and consistency.} Crowd-sourced speech corpora, while invaluable for bootstrapping low-resource ASR, introduce variability in recording conditions, microphone quality, background noise levels, and speaker characteristics that would be controlled in professionally recorded datasets. The Mozilla Common Voice platform, which provides the Kalenjin data used in this work, relies on volunteer contributors recording utterances through web browsers or mobile applications in uncontrolled acoustic environments. While community validation mechanisms help filter egregiously poor recordings, the resulting corpus inevitably contains heterogeneous acoustic conditions that the preprocessing pipeline must address.

\textbf{Linguistic resource gaps.} Beyond speech data, effective ASR systems benefit from pronunciation dictionaries that map words to phoneme sequences, morphological analyzers that decompose complex words into constituent morphemes, and large text corpora for training language models. For Kalenjin, none of these resources exist at scale. The absence of a pronunciation dictionary means that the model must learn grapheme-to-phoneme relationships implicitly from the speech-text pairs in the training data. The lack of large text corpora limits the effectiveness of language models in correcting acoustic model errors during decoding.

\textbf{Typological distance from pre-training languages.} While cross-lingual transfer learning mitigates data scarcity, its effectiveness depends partly on the typological similarity between the target language and the languages represented in the pre-training data. Kalenjin's tonal properties, agglutinative morphology, and Nilotic phonological inventory may be underrepresented in the pre-training distribution of XLS-R, which, despite its multilingual scope, is weighted toward Indo-European and other well-documented language families. This typological distance may limit the transferability of pre-trained representations and necessitate more extensive fine-tuning to adapt the model to Kalenjin-specific acoustic patterns.

\textbf{Training pathologies in low-resource settings.} Fine-tuning large pre-trained models on very small datasets introduces specific optimization challenges. The model may converge to degenerate solutions---such as predicting only blank tokens under CTC loss---if the learning rate, batch size, or training duration are not carefully calibrated. The risk of overfitting is elevated when the number of trainable parameters vastly exceeds the number of training examples, requiring regularization strategies such as dropout, layer freezing, and data augmentation. These pathologies, which are rarely encountered in high-resource settings, demand systematic diagnostic analysis and iterative refinement of the training protocol.

\textbf{Evaluation challenges.} Standard ASR evaluation metrics---Word Error Rate (WER) and Character Error Rate (CER)---may not fully capture the performance characteristics relevant to under-resourced languages. For agglutinative languages like Kalenjin, a single morpheme error can produce a word-level error even when the majority of the word is correctly transcribed, potentially inflating WER relative to the model's actual phonetic accuracy. CER provides a more granular assessment but does not capture word-level semantic accuracy. The interpretation of error rates must therefore account for the morphological properties of the target language.

\subsection{Research Objectives}

This research addresses the technical and operational challenges of building an ASR system for Kalenjin through a systematic, multi-stage approach. The primary objective is to develop a functional speech recognition system that demonstrates the viability of transfer learning from multilingual pre-trained models for extremely low-resource Nilotic languages. Specifically, this work pursues the following objectives:

\begin{enumerate}[label=\textbf{\arabic*.}, leftmargin=2em]
    \item \textbf{Design and implement a comprehensive preprocessing pipeline} for Kalenjin speech data that integrates audio signal processing (resampling, voice activity detection, silence trimming, peak normalization), Kalenjin-specific text normalization (orthographic standardization, character filtering, CTC delimiter insertion), and multi-dimensional quality assessment (signal-to-noise ratio estimation, spectral feature analysis, duration filtering, text quality validation).

    \item \textbf{Investigate training dynamics in extremely low-resource settings} by systematically analyzing the convergence behavior, failure modes, and optimization pathologies that arise when fine-tuning a 300-million-parameter model on fewer than 20 hours of transcribed speech. This includes documenting the degenerate blank-prediction behavior observed under suboptimal training configurations and identifying the critical factors---data volume, learning rate, training duration, and batch size---that determine whether the model learns meaningful transcription patterns.

    \item \textbf{Develop and evaluate a two-stage fine-tuning strategy} that separates initial adaptation of the Transformer encoder and CTC head (with frozen feature encoder) from subsequent end-to-end refinement with feature encoder unfreezing. This staged approach aims to balance the preservation of pre-trained acoustic representations with the need for language-specific acoustic adaptation.

    \item \textbf{Assess the impact of language model integration} on word-level accuracy by constructing a 5-gram statistical language model from the training transcriptions and incorporating it into beam search decoding. This component addresses the lexical coherence limitations inherent in purely acoustic CTC decoding.

    \item \textbf{Establish a reproducible baseline} for Kalenjin ASR that can serve as a foundation for future research, providing documented preprocessing procedures, training configurations, and evaluation protocols that enable systematic comparison of alternative approaches.
\end{enumerate}

\subsection{Research Questions}

The following research questions guide the experimental design and analysis presented in this work:

\begin{enumerate}[label=\textbf{RQ\arabic*.}, leftmargin=3em]
    \item To what extent can cross-lingual transfer learning from Wav2Vec2-XLS-R-300M, pre-trained on 436,000 hours of multilingual speech, compensate for the extreme data scarcity of Kalenjin (approximately 20 hours of transcribed speech)?

    \item What are the critical training configuration parameters---including data volume, learning rate, number of epochs, and feature encoder freezing strategy---that determine whether fine-tuning produces meaningful transcriptions versus degenerate solutions in extremely low-resource settings?

    \item Does a two-stage training strategy, in which the convolutional feature encoder is first frozen and subsequently unfrozen, yield measurable improvements over single-stage fine-tuning for Kalenjin ASR?

    \item To what degree does the integration of a statistical language model trained on limited Kalenjin text data improve word-level transcription accuracy beyond the acoustic model alone?
\end{enumerate}

\subsection{Significance of the Study}

This research contributes to the broader effort of extending speech technology to the world's under-resourced languages, with implications that span technical, social, and academic dimensions.

From a \textbf{technical perspective}, this work provides empirical evidence on the effectiveness of cross-lingual transfer learning for a Nilotic language, a language family that has received minimal attention in the ASR literature. The systematic documentation of training pathologies---including the degenerate blank-prediction failure mode and its resolution through data scaling and hyperparameter adjustment---offers practical guidance for researchers working with other extremely low-resource languages. The two-stage training strategy and its demonstrated 55.3\% relative CER reduction from Stage~1 to Stage~2 provides a concrete methodology that can be adapted to similar settings.

From a \textbf{social perspective}, the development of Kalenjin ASR technology has the potential to enhance digital inclusion for millions of speakers who currently lack access to voice-enabled services in their native language. Practical applications include voice-based agricultural information systems for rural farmers, automated transcription of local language media broadcasts, educational tools that support mother-tongue literacy, and healthcare communication aids in regions where Kalenjin is the primary language of daily interaction. While the current system's error rates preclude deployment in high-stakes applications, the established baseline demonstrates feasibility and provides a foundation for continued improvement.

From an \textbf{academic perspective}, this work contributes to the growing body of literature on low-resource ASR by providing a detailed case study that documents not only the final system performance but also the iterative process of diagnosing failures, refining methodologies, and achieving incremental improvements. The progression from 100\% WER (degenerate baseline) through 72.11\% (initial GPU training), 69.13\% (two-stage refinement), to 61.75\% (with language model integration) illustrates the cumulative impact of methodological decisions and provides a roadmap for researchers approaching similar challenges.

\subsection{Scope and Delimitations}

This study focuses specifically on read-speech recognition for Kalenjin using the Mozilla Common Voice v24.0 corpus. The following delimitations define the boundaries of this work:

\begin{itemize}[leftmargin=2em]
    \item The ASR system is trained and evaluated exclusively on read speech recorded through the Common Voice platform. Spontaneous conversational speech, code-switched utterances, and speech in noisy real-world environments are not addressed.

    \item The study treats Kalenjin as a single language entity without explicit modeling of sub-variety differences. While the training data likely contains contributions from speakers of multiple Kalenjin sub-varieties, no dialect identification or dialect-specific adaptation is performed.

    \item The language model is trained solely on the transcriptions from the training partition of the Common Voice corpus (11,057 sentences). No external Kalenjin text data from web sources, published texts, or other corpora is incorporated.

    \item Model training is conducted using a single pre-trained foundation model (Wav2Vec2-XLS-R-300M). Comparison with alternative foundation models (e.g., Whisper, HuBERT, WavLM) or alternative model sizes is deferred to future work.

    \item The evaluation is performed on the standard Common Voice test partition. No evaluation on external test sets or in-the-wild recordings is conducted.
\end{itemize}

\subsection{Organization of the Thesis}

The remainder of this thesis is organized as follows. Chapter~2 presents a comprehensive review of the literature on automatic speech recognition, self-supervised speech representation learning, low-resource ASR methodologies, and prior work on African language speech technology. Chapter~3 details the complete methodology, encompassing data acquisition, audio preprocessing, text normalization, quality assessment, model architecture, training configuration, and evaluation protocols. Chapter~4 presents the experimental results, including preprocessing statistics, training dynamics across multiple configurations, and final system performance with and without language model integration. Chapter~5 provides a detailed discussion of the results, analyzing error patterns, contextualizing performance within the broader low-resource ASR landscape, and identifying limitations. Chapter~6 concludes the thesis with a summary of contributions and directions for future research.

% ============================================================================
% REFERENCES
% ============================================================================
\bibliographystyle{plainnat}

\begin{thebibliography}{99}

\bibitem[Adelani et~al.(2022)]{adelani2022masakhaner}
Adelani, D.~I., Abbott, J., Neubig, G., D'souza, D., Kreutzer, J., Lignos, C., \ldots\ \& Buzaaba, H. (2022).
\newblock MasakhaNER 2.0: Africa-centric transfer learning for named entity recognition.
\newblock \textit{Proceedings of the 2022 Conference on Empirical Methods in Natural Language Processing}, 4488--4508.

\bibitem[Babu et~al.(2022)]{babu2022xls}
Babu, A., Wang, C., Tjandra, A., Lakhotia, K., Xu, Q., Goyal, N., \ldots\ \& Auli, M. (2022).
\newblock XLS-R: Self-supervised cross-lingual speech representation learning at scale.
\newblock \textit{arXiv preprint arXiv:2111.09296}.

\bibitem[Baevski et~al.(2020)]{baevski2020wav2vec}
Baevski, A., Zhou, Y., Mohamed, A., \& Auli, M. (2020).
\newblock wav2vec 2.0: A framework for self-supervised learning of speech representations.
\newblock \textit{Advances in Neural Information Processing Systems}, 33, 12449--12460.

\bibitem[Chan et~al.(2016)]{chan2016listen}
Chan, W., Jaitly, N., Le, Q., \& Vinyals, O. (2016).
\newblock Listen, attend and spell: A neural network that learns to read.
\newblock \textit{Proceedings of the IEEE International Conference on Acoustics, Speech and Signal Processing (ICASSP)}, 4960--4964.

\bibitem[Graves(2006)]{graves2006connectionist}
Graves, A., Fern{\'a}ndez, S., Gomez, F., \& Schmidhuber, J. (2006).
\newblock Connectionist temporal classification: Labelling unsegmented sequence data with recurrent neural networks.
\newblock \textit{Proceedings of the 23rd International Conference on Machine Learning}, 369--376.

\bibitem[Graves(2012)]{graves2012sequence}
Graves, A. (2012).
\newblock Sequence transduction with recurrent neural networks.
\newblock \textit{arXiv preprint arXiv:1211.3711}.

\bibitem[Hinton et~al.(2012)]{hinton2012deep}
Hinton, G., Deng, L., Yu, D., Dahl, G.~E., Mohamed, A., Jaitly, N., \ldots\ \& Kingsbury, B. (2012).
\newblock Deep neural networks for acoustic modeling in speech recognition: The shared views of four research groups.
\newblock \textit{IEEE Signal Processing Magazine}, 29(6), 82--97.

\bibitem[Joshi et~al.(2020)]{joshi2020state}
Joshi, P., Santy, S., Buber, A., Bali, K., \& Choudhury, M. (2020).
\newblock The state and fate of linguistic diversity and inclusion in the NLP world.
\newblock \textit{Proceedings of the 58th Annual Meeting of the Association for Computational Linguistics}, 6282--6293.

\bibitem[Magueresse et~al.(2020)]{magueresse2020low}
Magueresse, A., Carles, V., \& Heetderks, E. (2020).
\newblock Low-resource languages: A review of past work and future challenges.
\newblock \textit{arXiv preprint arXiv:2006.07264}.

\bibitem[Rabiner(1989)]{rabiner1989tutorial}
Rabiner, L.~R. (1989).
\newblock A tutorial on hidden Markov models and selected applications in speech recognition.
\newblock \textit{Proceedings of the IEEE}, 77(2), 257--286.

\bibitem[Radford et~al.(2023)]{radford2023robust}
Radford, A., Kim, J.~W., Xu, T., Brockman, G., McLeavey, C., \& Sutskever, I. (2023).
\newblock Robust speech recognition via large-scale weak supervision.
\newblock \textit{Proceedings of the 40th International Conference on Machine Learning}, 28492--28518.

\bibitem[Rottland(1982)]{rottland1982southern}
Rottland, F. (1982).
\newblock \textit{Die S{\"u}dnilotischen Sprachen: Beschreibung, Vergleichung und Rekonstruktion}.
\newblock Dietrich Reimer Verlag.

\bibitem[Toweett(1979)]{toweett1979study}
Toweett, T. (1979).
\newblock \textit{A Study of Kalenjin Linguistics}.
\newblock Kenya Literature Bureau.

\end{thebibliography}

\end{document}
